\documentclass[aspectratio=169, 11pt]{beamer}
\usepackage{beamer_preamble}

\begin{document}

% ==== Title ===============================================================
\titleslide{Lesson 4-3 \textperiodcentered\ Period 1}{Rational Expressions --- Equivalent \& Simplify}

% ==== Do Now ==============================================================
\begin{frame}{Do Now --- Explore \& Reason}{Algebra 2 \textperiodcentered\ Lesson 4-3 \textperiodcentered\ Period 1}
\phasebadges{1}{5}{notice \& wonder}
\vspace{0.2cm}
\begin{projcard}{Graph of \(y = x + 2\)}{slidegold}
Consider the graph of the function \(y = x + 2\).

\vspace{0.25cm}
\begin{enumerate}
  \setlength{\itemsep}{0.2cm}
  \item What is the domain of this function?
  \item Sketch a version that \textbf{excludes} \(x = 2\) from the domain.
  \item What kind of function might have a graph like that? Explain.
\end{enumerate}

\vspace{0.1cm}
\small Write 1 \textbf{notice} + 1 \textbf{wonder} on your packet. One sentence: what is the equation doing at that skipped \(x\)-value?
\end{projcard}
\end{frame}

% ==== Launch ==============================================================
\begin{frame}{Launch --- Equivalent + Simplify Rules}{Algebra 2 \textperiodcentered\ Lesson 4-3 \textperiodcentered\ Period 1}
\phasebadges{1--2}{12}{teacher models}
\vspace{0.2cm}
\begin{projcard}{Equivalent + Simplify Rules (keep on your page)}{slidegreen}
\begin{enumerate}
  \setlength{\itemsep}{0pt}
  \item \textbf{Factor} the numerator AND denominator completely.
  \item Cancel only common \textbf{FACTORS} --- never terms across \(+\) or \(-\).
  \item State domain restrictions: every zero of the \textbf{ORIGINAL} denominator is excluded --- including factors you canceled.
\end{enumerate}

\vspace{0.1cm}
\small
\textbf{Common Error:} Do NOT cancel terms that are added or subtracted across a fraction bar.
Factor completely FIRST, then cancel only identical factor pairs.

\vspace{0.1cm}
\textcolor{slideaccent}{\textbf{Key insight:} Canceling a factor never erases the domain restriction it imposed.}
\end{projcard}
\end{frame}

% ==== Explore =============================================================
\begin{frame}{Explore --- Think-Pair-Share (33 min)}{Algebra 2 \textperiodcentered\ Lesson 4-3 \textperiodcentered\ Period 1}
\phasebadges{1--2}{33}{student work}
\small\textcolor{slidegray}{7 items in order. Plan \(\sim\)4--5 min per item. Cut \#19 if pace slows.}

\vspace{0.1cm}
\begin{projcard}{Try It 1 \textperiodcentered\ Equivalent expression}{slideblue}
Write an expression equivalent to \((x-4)/x\) over domain \(\{x \mid x \ne 0 \text{ or } {-2}\}\).
\end{projcard}
\vspace{-0.15cm}
\begin{projcard}{Practice \#20 \textperiodcentered\ Equivalent over restricted domain}{slideblue}
Write an equivalent expression: \(x(x-2)/(x-5)\) for all \(x\) except \(5\) and \(-6\).
\end{projcard}
\vspace{-0.15cm}
\begin{projcard}{Practice \#22 \textperiodcentered\ Simplify + domain}{slideblue}
Simplify \((y^2-5y-24)/(y^2+3y)\). State the domain.
\end{projcard}
\vspace{-0.15cm}
\begin{projcard}{Practice \#24 \textperiodcentered\ Simplify + domain}{slideblue}
Simplify \((x^2+8x+15)/(x^2-x-12)\). State the domain.
\end{projcard}
\end{frame}

\begin{frame}{Explore (continued)}{Algebra 2 \textperiodcentered\ Lesson 4-3 \textperiodcentered\ Period 1}
\phasebadges{1--2}{33}{student work}
\vspace{0.2cm}
\begin{projcard}{Try It 2 \textperiodcentered\ Simplify by factoring and canceling}{slideblue}
Simplify each expression and state the domain.\\[4pt]
\textbf{a.}\ \((x^2+2x+1)/(x^3-2x^2-3x)\)
\hfill
\textbf{b.}\ \((x^3+4x^2-x-4)/(x^2+3x-4)\)
\end{projcard}
\vspace{0.1cm}
\begin{projcard}{Practice \#17 \textperiodcentered\ Construct Arguments}{slideblue}
Explain how you can tell whether a rational expression is in simplest form.
\end{projcard}
\vspace{0.1cm}
\begin{projcard}{Practice \#19 \textperiodcentered\ Domain restriction (extension)}{slideaccent}
If the denominator is \(x^3+3x^2-10x\), what value(s) must be restricted from the domain?
\end{projcard}
\end{frame}

% ==== Share / Summary =====================================================
\begin{frame}{Share / Summary}{Algebra 2 \textperiodcentered\ Lesson 4-3 \textperiodcentered\ Period 1}
\phasebadges{2}{5}{academic conversation}
\vspace{0.2cm}
\begin{projcard}{Sentence frame \textperiodcentered\ Domain \& Simplify}{slideblue}
``The factored form is \underline{\hspace{1.8cm}}, so the domain excludes \(x =\) \underline{\hspace{1cm}}.
After canceling the factor \underline{\hspace{1.5cm}}, the simplified expression is \underline{\hspace{1.5cm}}
with the \textbf{same} domain restriction.''
\end{projcard}

\vspace{0.15cm}
\begin{projcard}{Bridge to Period 2}{slideaccent}
Next class: multiply and divide rational expressions.\\
The same factor-and-cancel logic applies --- plus you'll work with products and quotients.\\
\textbf{Key rule carries over:} always state domain restrictions before simplifying.
\end{projcard}
\end{frame}

% ==== Exit Ticket =========================================================
\begin{frame}{Exit Ticket --- Summary Recap}{Algebra 2 \textperiodcentered\ Lesson 4-3 \textperiodcentered\ Period 1}
\phasebadges{1--2}{2}{summary recap --- not graded}
\vspace{0.2cm}
\begin{projcard}{Complete on your exit ticket}{slidegold}
To simplify \((x^2+3x-10)/(x^2-4)\) I first factor to \underline{\hspace{3.5cm}}.

\vspace{0.3cm}
The restricted domain values are \(x \ne\) \underline{\hspace{2cm}}.

\vspace{0.3cm}
After canceling the factor \underline{\hspace{2cm}}, my simplified expression is \underline{\hspace{2cm}}.

\vspace{0.3cm}
One biggest thing I learned today: \underline{\hspace{5cm}}.
\end{projcard}
\end{frame}

\end{document}
